\chapter{Monopoly and Pricing}\label{ch:b1:monopoly-pricing}

An airline sells seats on the same flight at five prices. A cinema charges
students less than other adults for the same film. A phone plan has a
monthly fee and a price per gigabyte. None of this happens in the
competitive markets of the last chapters, where every unit sells at one
price equal to marginal cost. It happens when a seller faces a
downward-sloping demand and can choose its price, and the way it chooses is
the subject of this chapter.

Market power is hard to see directly, because marginal cost is not
published. One margin that is published is the gap between the rate banks
charge on loans and the rate they pay on deposits.

\begin{openingchart}{Banks' lending rate minus deposit rate against GDP per
person at purchasing-power parity, 110 countries, 2019, with the fitted
line.}{World Bank, World Development Indicators (FR.INR.LNDP,
NY.GDP.PCAP.PP.KD)}
\begin{tikzpicture}
\begin{axis}[omscatter, xmode=log, log ticks with fixed point, xtickten={3,4,5}, xmin=700, xmax=150000, ymin=0, ymax=40,
  xlabel={GDP per person, PPP (international dollars, log scale)}, ylabel={Lending minus deposit rate (points)},
  width=0.74\linewidth, height=0.44\linewidth]
  \addplot[only marks, mark=*, mark size=1.3pt, omDef, fill opacity=0.6] table[x=gdp, y=spread, col sep=comma] {figdata/bachelor-1/15-monopoly-pricing/out/spreads_2019.csv};
  \addplot[omThm, thick, sharp plot, mark=none] table[x=gdp, y=fit, col sep=comma] {figdata/bachelor-1/15-monopoly-pricing/out/spreads_fit.csv};
  \addplot[only marks, mark=none, nodes near coords, point meta=explicit symbolic, nodes near coords style={font=\scriptsize, anchor=west}]
    table[x=gdp, y=spread, meta=country, col sep=comma] {figdata/bachelor-1/15-monopoly-pricing/out/spreads_points.csv};
\end{axis}
\end{tikzpicture}
\end{openingchart}

The median spread is $5.88$ points; a tenth of the countries have less than
$2.29$ and a tenth more than $12.91$. Brazil's spread is $32.04$ points and
Madagascar's $36$, against $2.98$ in Switzerland. % figdata: spreads_stats, spreads_points
Spreads are narrower in richer countries: the median is $7.73$ points in the
poorer half and $4.90$ in the richer half, and the fitted line falls by $5.61$
points for each tenfold rise in income. % figdata: spreads_stats
A spread is not a \omterm{def:b1:monopoly-pricing:markup}{markup}: it also pays for the risk that loans are not
repaid, for reserves the bank must hold and for the cost of running
branches. But part of it is the price of market power in banking, and the
tools of this chapter say what part that could be.

\begin{recall}
A monopoly, marginal revenue, \omterm{def:b1:monopoly-pricing:natural-monopoly}{natural monopoly} and charging different buyers
different prices were met at school. The price elasticity of demand is in
\cref{ch:b1:comparative-statics}; surplus as an integral and the
\omterm{def:b1:partial-equilibrium-welfare:dwl}{deadweight loss} in \cref{ch:b1:partial-equilibrium-welfare}.
\end{recall}

\section{Monopoly with calculus}\label{sec:b1:monopoly-pricing:monopoly}

A single seller faces the inverse demand $p(q)$ and the cost $C(q)$. Its
profit is $\pi(q) = p(q)\,q - C(q)$, and at a maximum
\[
MR(q) = p(q) + q\,p'(q) = C'(q) = MC(q).
\]
Write $\varepsilon = \frac{\dd q}{\dd p}\frac pq < 0$ for the elasticity of demand. Then $qp'(q) = p/\varepsilon$, and
\[
MR = p\left(1 + \frac1\varepsilon\right).
\]

\begin{definition}[Lerner index]\label{def:b1:monopoly-pricing:lerner}
\index{Lerner index}
The \emph{Lerner index} of a firm is $L = (p - MC)/p$, the share of its price
above marginal cost.
\end{definition}

\begin{definition}[Markup]\label{def:b1:monopoly-pricing:markup}
\index{markup}
The \emph{markup} of a firm is the ratio of its price to its marginal cost,
$\mu = p/MC$. A competitive firm has $\mu = 1$ and $L = 0$; and $L = 1 - 1/\mu$.
\end{definition}

\begin{proposition}[The inverse-elasticity rule]\label{prop:b1:monopoly-pricing:inverse-elasticity}
At the profit-maximising price of a monopolist, $L = -1/\varepsilon$, so
$\mu = \varepsilon/(1 + \varepsilon)$. A monopolist with positive marginal cost never sets a price at
which demand is inelastic, $|\varepsilon| \le 1$.
\end{proposition}

\begin{proof}
$MR = MC$ gives $p(1 + 1/\varepsilon) = MC$, so $(p - MC)/p = -1/\varepsilon$. Where
$|\varepsilon| \le 1$, $MR = p(1 + 1/\varepsilon) \le 0 < MC$: cutting output raises revenue and lowers
cost, so profit rises.
\end{proof}

The rule says that market power is limited by the buyers' ability to walk
away. A seller facing an elasticity of $-2$ prices at twice marginal cost;
one facing $-11$, at $1.1$ times. And it gives a \omterm{def:b1:partial-equilibrium-welfare:dwl}{deadweight loss}: the units
between the monopoly quantity and the competitive one, each worth more to
some buyer than it costs to make, are not sold.

\begin{omfigure}
\begin{tikzpicture}
\begin{axis}[omrecord, grid=none, xlabel={Quantity $q$}, ylabel={Price}, xmin=0, xmax=105, ymin=0, ymax=105,
  width=0.52\linewidth, height=0.42\linewidth]
  \fill[omProp!35] (axis cs:40,60) -- (axis cs:80,20) -- (axis cs:40,20) -- cycle;
  \addplot[demand, domain=0:100] {100 - x};
  \addplot[shifted, domain=0:50] {100 - 2*x};
  \addplot[supply, domain=0:105] {20};
  \addplot[guide] coordinates {(0,60) (40,60) (40,0)};
  \node[font=\footnotesize, anchor=west] at (axis cs:86,24) {$MC$};
  \node[font=\footnotesize, anchor=west] at (axis cs:44,14) {$MR$};
  \node[font=\footnotesize, anchor=west] at (axis cs:70,38) {$D$};
  \node[font=\footnotesize] at (axis cs:52,32) {DWL};
\end{axis}
\end{tikzpicture}

{\small Demand $p = 100 - q$, marginal revenue $100 - 2q$ and marginal cost 20.
The monopolist sells 40 at 60, against 80 at 20 in competition; the
triangle, 800, is the \omterm{def:b1:partial-equilibrium-welfare:dwl}{deadweight loss}.}
\end{omfigure}

\section{Price discrimination}\label{sec:b1:monopoly-pricing:discrimination}

\begin{definition}[Price discrimination]\label{def:b1:monopoly-pricing:discrimination}
\index{price discrimination}
A seller \emph{price-discriminates} when it sells units of the same good at
prices whose differences are not explained by differences in cost. It is
of the \emph{first degree} when each unit sells at the most its buyer would
pay; of the \emph{second degree} when all buyers face the same menu of
prices that depend on the quantity or version bought, and choose from it;
of the \emph{third degree} when the seller sets a different price for each
observable group of buyers.
\end{definition}

Discrimination needs two things: the seller must be able to tell buyers
apart, by group or by what they choose, and the buyers who pay less must
not be able to resell to those who pay more. Services, which cannot be
resold, lend themselves to it.

\emph{First degree.} If the seller charges each buyer exactly his
willingness to pay, it sells every unit worth more than its marginal cost:
the quantity is the competitive one, there is no \omterm{def:b1:partial-equilibrium-welfare:dwl}{deadweight loss}, and the
whole surplus goes to the seller.

\emph{Third degree.} With two groups and a common marginal cost, the seller
chooses $q_1$ and $q_2$ to maximise $p_1(q_1)q_1 + p_2(q_2)q_2 - C(q_1 + q_2)$, and so
$MR_1 = MR_2 = MC$.

\begin{proposition}[The more elastic group pays less]\label{prop:b1:monopoly-pricing:third}
Under third-degree discrimination, $p_1(1 + 1/\varepsilon_1) = p_2(1 + 1/\varepsilon_2)$, so the group whose
demand is more elastic at its price pays the lower price.
\end{proposition}

\begin{proof}
$MR_i = p_i(1 + 1/\varepsilon_i)$ and both equal $MC$. If $|\varepsilon_1| > |\varepsilon_2|$, then
$1 + 1/\varepsilon_1 > 1 + 1/\varepsilon_2$, so $p_1 < p_2$.
\end{proof}

Students and pensioners get cheaper tickets because they are more price
sensitive, not because they cost less to serve. The welfare effect is
ambiguous. If discrimination opens a market the single price left unserved,
it adds surplus. With linear demands and both groups served, total output
is the same as under the single price but is allocated less well, and
\omterm{def:b1:partial-equilibrium-welfare:surplus}{total surplus} falls (\cref{exo:b1:monopoly-pricing:welfare}).

\emph{Second degree.} When the seller cannot tell buyers apart, it can let
them sort themselves: a first-class and an economy seat, a quantity
discount, a ticket that must be bought weeks ahead. How to design such a
menu, so that each type of buyer picks the option meant for it, is the
screening problem of Year~2.

\begin{record}[Price discrimination in a Broadway theatre]
A study of the ticket prices of one Broadway play estimated buyers'
demands and computed what would happen under a single price. The
observed discrimination --- different prices for different seats, and
discounts for some buyers --- raised the producer's profit by about $5\%$
relative to the best uniform price, while total consumer welfare was almost
the same under both (Leslie 2004). % ledger: b1:leslie2004
\end{record}

\section{Two-part tariffs and bundling}\label{sec:b1:monopoly-pricing:tariffs}

\begin{definition}[Two-part tariff]\label{def:b1:monopoly-pricing:two-part}
\index{two-part tariff}
A \emph{two-part tariff} charges a fixed fee for the right to buy plus a
price per unit bought.
\end{definition}

\begin{proposition}[The optimal two-part tariff with identical buyers]\label{prop:b1:monopoly-pricing:two-part}
If all buyers have the same demand and marginal cost is constant, the most
profitable \omterm{def:b1:monopoly-pricing:two-part}{two-part tariff} sets the unit price equal to marginal cost and
the fee equal to each buyer's \omterm{def:b1:consumer-welfare:consumer-surplus}{consumer surplus} at that price.
\end{proposition}

\begin{proof}
At a unit price $p$, a buyer accepts any fee up to his \omterm{def:b1:consumer-welfare:consumer-surplus}{consumer surplus}
$\mathrm{CS}(p)$, so the seller earns $\mathrm{CS}(p) + (p - c)q(p)$ per buyer: the buyer's
surplus plus the profit on units, that is the \omterm{def:b1:partial-equilibrium-welfare:surplus}{total surplus} of the units
traded. This is largest at $p = c$ (\cref{prop:b1:partial-equilibrium-welfare:efficiency}).
\end{proof}

With identical buyers the \omterm{def:b1:monopoly-pricing:two-part}{two-part tariff} achieves first-degree
discrimination: no \omterm{def:b1:partial-equilibrium-welfare:dwl}{deadweight loss}, and all the surplus to the seller.
With different buyers, a single fee either excludes those with small
demands or leaves surplus to those with large ones, and the best unit price
is above marginal cost.

\begin{omfigure}
\begin{tikzpicture}
\begin{axis}[omrecord, grid=none, xlabel={Trips a year}, ylabel={Dalens a trip}, xmin=0, xmax=42, ymin=0, ymax=42,
  width=0.48\linewidth, height=0.4\linewidth]
  \fill[omDef!22] (axis cs:0,10) -- (axis cs:0,40) -- (axis cs:30,10) -- cycle;
  \addplot[demand, domain=0:40] {40 - x};
  \addplot[supply, domain=0:42] {10};
  \node[font=\footnotesize] at (axis cs:9,18) {fee $= 450$};
  \node[font=\footnotesize, anchor=south west] at (axis cs:31,10) {$MC = 10$};
\end{axis}
\end{tikzpicture}

{\small A resident of Brenmoor wants $40 - p$ ferry trips a year (the weekend
problem). Charged 10 a trip, he makes 30 and would pay up to 450 a year for
the right to do so: the fee of the best \omterm{def:b1:monopoly-pricing:two-part}{two-part tariff}.}
\end{omfigure}

\begin{definition}[Bundling]\label{def:b1:monopoly-pricing:bundling}
\index{bundling}
\emph{Bundling} is selling two or more goods together as a package at one
price, alone or alongside the goods sold separately.
\end{definition}

\omterm{def:b1:monopoly-pricing:bundling}{Bundling} pays when buyers' values for the goods are negatively correlated.
If one buyer values a newspaper at 8 and its crossword app at 3, and another
values them at 3 and 8, separate prices earn at most 16, but the bundle at 11
sells to both and earns 22. The bundle evens out buyers' values, which
makes demand for it less elastic, so it acts like \omterm{def:b1:monopoly-pricing:discrimination}{price discrimination}
(\cref{exo:b1:monopoly-pricing:bundle}).

\section{Natural monopoly}\label{sec:b1:monopoly-pricing:natural}

\begin{definition}[Natural monopoly]\label{def:b1:monopoly-pricing:natural-monopoly}
\index{natural monopoly}\index{subadditive cost}
An industry is a \emph{natural monopoly} when its \omterm{def:b1:cost-minimisation:cost-function}{cost function} is
\emph{subadditive} over the relevant outputs: one firm can produce any total
output more cheaply than two or more firms sharing it,
$C(q_1 + q_2) < C(q_1) + C(q_2)$.
\end{definition}

A large fixed cost and a constant marginal cost, $C = F + cq$, are the usual
case: a water network, a railway track, an electricity grid. Two firms pay
the fixed cost twice. Average cost $F/q + c$ falls with output and stays above
marginal cost, so pricing at marginal cost --- the efficient price --- loses
$F$. A regulator who will not pay a subsidy must let the price exceed
marginal cost.

\begin{definition}[Average-cost pricing and price-cap regulation]\label{def:b1:monopoly-pricing:regulation}
\index{average-cost pricing}\index{price-cap regulation}\index{RPI$-X$}
Under \emph{average-cost pricing} the regulator sets the price at which
revenue just covers total cost. Under \emph{price-cap regulation} it sets a
ceiling on the firm's prices that rises each year by the rate of inflation
minus a number $X$ fixed in advance for several years (RPI$-X$, after the
retail price index).
\end{definition}

The two differ in incentives. A firm whose price is reset to its average
cost gains nothing by cutting cost, since the price follows the cost down.
A firm under a cap keeps whatever it saves until the cap is next reset.

\begin{omfigure}
\begin{tikzpicture}
\begin{axis}[omrecord, grid=none, xlabel={Trips (thousands a year)}, ylabel={Dalens a trip}, xmin=0, xmax=200, ymin=0, ymax=80,
  width=0.56\linewidth, height=0.42\linewidth, samples=80]
  \fill[omProp!30] (axis cs:0,10) rectangle (axis cs:190,17.89);
  \addplot[demand, domain=40:200] {48 - x/5};
  \addplot[demand, domain=0:40] {80 - x};
  \addplot[supply, domain=0:200] {10};
  \addplot[shifted, domain=31:200] {10 + 1500/x};
  \addplot[guide] coordinates {(0,21.19) (134.05,21.19) (134.05,0)};
  \addplot[guide] coordinates {(190,0) (190,17.89)};
  \node[font=\footnotesize, anchor=west] at (axis cs:38,53) {$AC$};
  \node[font=\footnotesize, anchor=south west] at (axis cs:2,10) {$MC$};
  \node[font=\footnotesize, anchor=west] at (axis cs:12,70) {$D$};
  \node[font=\footnotesize] at (axis cs:95,14) {loss at $p = MC$};
\end{axis}
\end{tikzpicture}

{\small The ferry of the weekend problem: a fixed cost of 1,500 thousand
dalens a year and a marginal cost of 10 a trip. \omterm{def:b1:consumer-choice:market-demand}{Market demand} is
$240 - 5p$ for prices up to 40. At $p = 10$ the ferry carries 190 thousand trips
and loses 1,500; the average-cost price, 21.19, covers its costs with 134
thousand trips.}
\end{omfigure}

\begin{record}[Price caps for utilities]
The price cap was proposed in a 1983 report on how to regulate British
Telecom, then about to be privatised, and applied to it from 1984; the
caps were reviewed every five years. It was then used for gas,
electricity, water, railways and airports (Stern 2003). % ledger: b1:rpi-x
The model's case for it was the incentive to cut costs between reviews.
What happened is that the formulas grew complex --- water's included an
extra term for investment --- and that each review reset $X$ with an eye on
the firm's profits, so a firm that cut costs could expect part of its gain
to be taken back at the next review: the incentive works only between
reviews, and only as far as firms trust how the next one will be done
(Stern 2003). % ledger: b1:rpi-x
\end{record}

\section{Measuring market power}\label{sec:b1:monopoly-pricing:measuring}

The \omterm{def:b1:monopoly-pricing:lerner}{Lerner index} needs marginal cost, which accounts do not show. There are
two ways round it. One estimates the elasticity of demand a firm faces and
applies the inverse-elasticity rule; this needs a model of demand. The
other uses production: a cost-minimising firm pays for a flexible input up
to the point where its output elasticity equals its spending on the input
divided by marginal cost times output; so the ratio of that elasticity to
the input's share of \emph{revenue} is price over marginal cost, the
\omterm{def:b1:monopoly-pricing:markup}{markup}. This needs an estimate of the \omterm{def:b1:technology:production-function}{production function} but no
demand.

\begin{model}{Monopoly pricing}
\assumes One seller who knows its demand and its costs, and no entry.
\predicts A price above marginal cost, with the inverse-elasticity \omterm{def:b1:monopoly-pricing:markup}{markup};
higher prices to less elastic groups when they can be told apart; a
\omterm{def:b1:partial-equilibrium-welfare:dwl}{deadweight loss} that discrimination may shrink or enlarge.
\tested \omterm{def:b1:monopoly-pricing:discrimination}{Price discrimination} in a Broadway theatre raised profit by about
$5\%$ (Leslie 2004); % ledger: b1:leslie2004
\omterm{def:b1:monopoly-pricing:markup}{markups} measured from firms' accounts rose in the United States after 1980
(De Loecker, Eeckhout and Unger 2020), a finding whose size is debated.
% ledger: b1:deloecker2020
\breaks With entry and new products (Year~3); with durable goods, where the
monopolist competes with its own future sales; with regulated firms.
\end{model}

\begin{debate}[Have markups risen?]
Applying the production approach to the accounts of listed US firms since
1955, one study finds that the aggregate \omterm{def:b1:monopoly-pricing:markup}{markup} rose from $21\%$ above
marginal cost in 1980 to $61\%$ by the time of writing, with the median firm's \omterm{def:b1:monopoly-pricing:markup}{markup} unchanged and the
rise coming from the top of the distribution and from the growth of
\omterm{def:b1:monopoly-pricing:markup}{high-markup} firms; the average profit rate rose from $1\%$ to $8\%$ (De Loecker,
Eeckhout and Unger 2020). % ledger: b1:deloecker2020
Its weakness: listed firms are not the whole economy, and if the costs
treated as variable leave out overhead costs that have grown, a shift
towards overhead looks like a rise in \omterm{def:b1:monopoly-pricing:markup}{markups}. A review of three ways of measuring \omterm{def:b1:monopoly-pricing:markup}{markups} finds
that all rest on cost minimisation, that their answers differ widely, and
that existing methods cannot tell whether \omterm{def:b1:monopoly-pricing:markup}{markups} have been stable or have
risen modestly (Basu 2019). % ledger: b1:basu2019
Its weakness: it does not settle the question either. What would settle it
is direct evidence on marginal costs in many industries.
\end{debate}

\section*{Exercises}\label{sec:b1:monopoly-pricing:exercises}

\begin{exercise}[$\star$]\label{exo:b1:monopoly-pricing:basic}
A monopolist faces $p = 100 - q$ with $MC = 20$. Find the price, the quantity,
the profit and the \omterm{def:b1:partial-equilibrium-welfare:dwl}{deadweight loss}.
\end{exercise}

\begin{exercise}[$\star$]\label{exo:b1:monopoly-pricing:lerner-ex}
A firm faces an elasticity of $-4$ at its chosen price. Compute its \omterm{def:b1:monopoly-pricing:lerner}{Lerner
index} and its \omterm{def:b1:monopoly-pricing:markup}{markup}.
\end{exercise}

\begin{exercise}[$\star$]\label{exo:b1:monopoly-pricing:read}
From the opening chart's table, the spread was $32.04$ points in Brazil and
$2.98$ in Switzerland in 2019. % figdata: spreads_points
List three reasons other than market power that could explain the gap.
\end{exercise}

\begin{exercise}[$\star$]\label{exo:b1:monopoly-pricing:predict}
A museum can tell residents from tourists. Residents' demand is more
elastic at any price. Predict which group pays more.
\end{exercise}

\begin{exercise}[$\star\star$]\label{exo:b1:monopoly-pricing:third-ex}
A firm with $MC = 20$ sells in two markets with demands $q_1 = 100 - p_1$ and
$q_2 = 60 - p_2$. Find the discriminating prices and the single price, with the
quantities.
\end{exercise}

\begin{exercise}[$\star\star$]\label{exo:b1:monopoly-pricing:welfare}
With the numbers of \cref{exo:b1:monopoly-pricing:third-ex}, compute \omterm{def:b1:partial-equilibrium-welfare:surplus}{total
surplus} and the \omterm{def:b1:partial-equilibrium-welfare:dwl}{deadweight loss} under discrimination and under the single
price.
\end{exercise}

\begin{exercise}[$\star\star$]\label{exo:b1:monopoly-pricing:tpt}
Each buyer demands $q = 10 - p$ and $MC = 2$. Find the best \omterm{def:b1:monopoly-pricing:two-part}{two-part tariff}.
\end{exercise}

\begin{exercise}[$\star\star$]\label{exo:b1:monopoly-pricing:bundle}
Two buyers value a newspaper and an app at $(8, 3)$ and $(3, 8)$; costs are
zero. Compare the best separate prices with the best bundle price. What
would happen with values $(8, 8)$ and $(3, 3)$?
\end{exercise}

\begin{exercise}[$\star\star$]\label{exo:b1:monopoly-pricing:cap}
A regulator moves a water company from \omterm{def:b1:monopoly-pricing:regulation}{average-cost pricing} to a price cap
with $X = 3\%$, when inflation is $2\%$. Predict the path of the price over the
next five years and what the company does about its costs.
\end{exercise}

\begin{exercise}[$\star\star\star$]\label{exo:b1:monopoly-pricing:derive}
Derive $MR = p(1 + 1/\varepsilon)$ and the inverse-elasticity rule, and show that a
monopolist with positive marginal cost never prices where $|\varepsilon| < 1$.
\end{exercise}

\begin{exercise}[$\star\star\star$]\label{exo:b1:monopoly-pricing:ramsey}
A \omterm{def:b1:monopoly-pricing:natural-monopoly}{natural monopoly} has $C = 1200 + 20q$ and faces $q = 100 - p$. Find the lowest
price at which it breaks even, and the \omterm{def:b1:partial-equilibrium-welfare:dwl}{deadweight loss} there. Why is this
the best a regulator can do without a subsidy if the firm sells one good?
With several goods, which should bear the highest \omterm{def:b1:monopoly-pricing:markup}{markup}?
\end{exercise}

\begin{exercise}[$\star\star\star$]\label{exo:b1:monopoly-pricing:refute}
What observation would refute the claim that the measured rise in \omterm{def:b1:monopoly-pricing:markup}{markups}
since 1980 reflects a rise in market power rather than in fixed costs?
\end{exercise}

\section{Problem: The Ferry between Brenland and Kestria}\label{sec:b1:monopoly-pricing:problem}

\begin{problem}[{Weekend problem --- what a season pass removes}]\label{pb:b1:monopoly-pricing:ferry}
One operator runs the ferry between Brenmoor and Kestria (invented). A trip
costs it 10 dalens at the margin. There are 4,000 residents, each wanting
$40 - p$ trips a year, so residents' demand is $160 - 4p$ thousand trips;
tourists' demand is $80 - p$ thousand trips. Money below is in thousands of
dalens a year.

\medskip\noindent\textbf{Part I --- A single price.}
\begin{enumerate}
  \item Write total demand for prices up to 40.
  \item Find the profit-maximising price and quantity.
  \item How many trips does each group make?
  \item Compute the profit, and check that serving tourists alone at their best price would earn less.
  \item Compute the \omterm{def:b1:monopoly-pricing:lerner}{Lerner index} and check it against the elasticity of total demand.
  \item Compute each group's \omterm{def:b1:consumer-welfare:consumer-surplus}{consumer surplus}.
  \item Compute the \omterm{def:b1:partial-equilibrium-welfare:dwl}{deadweight loss}.
\end{enumerate}

\medskip\noindent\textbf{Part II --- Residents' and tourists' prices.}
\begin{enumerate}[resume]
  \item Find the residents' price and trips.
  \item Find the tourists' price and trips.
  \item Check $MR_1 = MR_2 = MC$, and compare the two elasticities.
  \item Compute the profit.
  \item Compute each group's \omterm{def:b1:consumer-welfare:consumer-surplus}{consumer surplus}.
  \item Compute the \omterm{def:b1:partial-equilibrium-welfare:dwl}{deadweight loss} and compare with Part I.
  \item Explain why total trips do not change and why surplus falls.
\end{enumerate}

\medskip\noindent\textbf{Part III --- A season pass for residents.} Tourists
still pay their price of Part II.
\begin{enumerate}[resume]
  \item What is a resident's \omterm{def:b1:consumer-welfare:consumer-surplus}{consumer surplus} at a price of 10 a trip?
  \item Design the most profitable pass: fee and price per trip. How many trips do residents make?
  \item Compute the profit from residents and in all.
  \item Compute the \omterm{def:b1:partial-equilibrium-welfare:dwl}{deadweight loss}.
  \item Why can the operator not sell tourists a pass on the same terms?
  \item What share of the single-price \omterm{def:b1:partial-equilibrium-welfare:dwl}{deadweight loss} of Part I does the pass remove?
\end{enumerate}

\medskip\noindent\textbf{Part IV --- The regulator.} Running the ferry costs
1,500 a year whatever the number of trips; a single price is charged.
\begin{enumerate}[resume]
  \item Show that the ferry is a \omterm{def:b1:monopoly-pricing:natural-monopoly}{natural monopoly}.
  \item What happens at a price equal to marginal cost?
  \item Find the average-cost price and the trips.
  \item Compute the \omterm{def:b1:partial-equilibrium-welfare:dwl}{deadweight loss} at that price.
  \item The regulator sets a cap at the average-cost price, rising by inflation ($2\%$) minus $X = 3\%$. What is next year's cap?
  \item The operator cuts its marginal cost to 9. With both costs, 9 a trip and 1,500 a year, stated in next year's money, compute its profit next year under the cap, and under \omterm{def:b1:monopoly-pricing:regulation}{average-cost pricing}.
  \item State the result: the share of the single-price \omterm{def:b1:partial-equilibrium-welfare:dwl}{deadweight loss} that the residents' pass removes.
\end{enumerate}
\end{problem}
